\section{Degree, Trace, Pushforward, and Pullback}
\label{sec:degree-trace}

\begin{layman}
This section is about how Riemann surfaces \emph{interact} when
mapped to one another. A holomorphic map \(f : X \to Y\) between
compact Riemann surfaces is either constant or a \emph{finite
branched covering} — a map where every point of \(Y\) has the same
finite number of preimages on \(X\), with multiplicities at a
finite set of \emph{branch points}. The number is the
\emph{degree} of \(f\), and is the first invariant we attach to
the map.

Two natural operations come along for the ride. \emph{Pullback}
\(f^*\) takes a holomorphic 1-form on the target \(Y\) and pulls
it up to a 1-form on the source \(X\), via the chain rule: in
local coordinates, \(f^*(h(w)\,dw) = h(f(z))\,f'(z)\,dz\). It is
linear and well-behaved everywhere. \emph{Trace} \(\tr_f\), or
\emph{pushforward}, goes the opposite direction: takes a 1-form
on the source and pushes it down to one on the target. The
construction is delicate — one must sum over the local inverse
branches of \(f\) on the target, define a putative form there,
and then prove that the result extends holomorphically across
the branch points (where the inverse branches collide). The
extension uses symmetric-function tricks; the result is a
complex-linear map \(\tr_f : H^0(X, \Omega^1_X) \to
H^0(Y, \Omega^1_Y)\).

The combination of pullback and trace, when dualised, descends to
the Jacobian: the pullback induces a holomorphic group
homomorphism \(f^* : \Jac(Y) \to \Jac(X)\), and the trace induces
the dual \(f_* : \Jac(X) \to \Jac(Y)\). The pair satisfies the
expected functoriality: identity goes to identity, composition
respects composition. The fourth and final anti-hack theorem of
the challenge then asserts the \emph{trace identity}:
\(f_*(f^* P) = \deg(f) \cdot P\)
for every Jacobian element \(P\). Equivalently, ``pulling back
then pushing forward equals multiplication by the degree.'' This
is forced upon us by the underlying linear identity on 1-forms:
\(\tr_f \circ f^* = \deg(f)\).
The geometric mnemonic: pulling back creates \(d = \deg(f)\)
copies of every form; tracing them down sums those \(d\) copies
back together.

Why bother? The four challenge anti-hack theorems are all in
play by now: genus tied to topology (section 4), Abel--Jacobi
injective (section 8), Jacobian a complex Lie group (section 7),
and \emph{this} section's trace identity. Together they rule out
every cheap construction that might define push and pull as ad
hoc homomorphisms unrelated to the underlying geometry. With the
trace identity, push and pull are forced to behave as classical
geometers always wanted them to behave.
\end{layman}

Let \(f:X\to Y\) be holomorphic. If \(f\) is constant, define
\(\deg(f)=0\). Otherwise \(f\) is a finite branched covering, and
\(\deg(f)\) is the cardinality of a generic fiber counted with multiplicity.

\begin{lemma}[Pullback of holomorphic one-forms]
\label{lem:pullback-forms}
\uses{lem:pbf-r1}
Pullback gives a complex-linear map
\[
  f^{*}:\HY\to\HX.
\]
\lean{JacobianChallenge.TraceDegree.pullbackFormsLinearMap}
\leanok
\end{lemma}

\begin{layman}
Given a holomorphic map between compact Riemann surfaces, the
\emph{pullback} on 1-forms is the most natural way to transport a
1-form on the target back to a 1-form on the source: if you know how
to measure infinitesimal directions downstream, you can measure
infinitesimal directions upstream by first applying the map and then
measuring downstream. Locally, if a target 1-form is \(h(w)\,dw\), its
pullback through a map \(w = f(z)\) is just \(h(f(z))\,f'(z)\,dz\) —
the chain rule. This operation is complex-linear (sums and scalar
multiples come along for the ride) and lands in honest holomorphic
1-forms. Why bother? Pullback gives a well-behaved arrow on the spaces
of holomorphic 1-forms, in the opposite direction to the map itself.
Dualizing this arrow will produce the pullback on Jacobians, the first
half of the push/pull pair the trace identity needs.
\end{layman}

\begin{proof}
\leanok
In local coordinates, if \(\eta=h(w)\,dw\), then
\[
  f^{*}\eta = h(f(z))f'(z)\,dz,
\]
which is holomorphic. Linearity is immediate.
\end{proof}
\begin{lemma}[Trace of holomorphic one-forms]
\label{lem:trace-forms}
\uses{lem:trf-r1}
For nonconstant \(f\), there is a complex-linear trace map
\[
  \tr_f:\HX\to\HY.
\]
\lean{JacobianChallenge.TraceDegree.analyticPushforward}
\leanok
\end{lemma}

\begin{theorem}[Trace--pullback identity for holomorphic 1-forms]
\label{thm:trace-pullback-identity}
\uses{lem:trace-forms,lem:pullback-forms}
For nonconstant holomorphic \(f : X \to Y\) and \(\eta \in H^{0}(Y, \Omega^{1})\),
\[
  \tr_{f}(f^{*}\eta) = \deg(f) \cdot \eta.
\]
\lean{JacobianChallenge.HolomorphicForms.trace_pullback_identity}
\notready
\end{theorem}

\begin{layman}
Given a nonconstant holomorphic map between compact Riemann surfaces,
...
the \emph{trace} on 1-forms is a natural way to push a 1-form on the
source down to a 1-form on the target. Away from the (finite) branch
locus, the map is a local degree-\(d\) covering with \(d\) inverse
branches; given a source 1-form, evaluate it along each of the \(d\)
branches and add. The result is a target 1-form away from the branch
points, and a careful analysis (using that symmetric functions of the
branches extend holomorphically) shows the form continues smoothly
across the branch points to give a global holomorphic 1-form on the
target. The operation is complex-linear, just like pullback. Why
bother? Trace is the dual partner to pullback. In the Jacobian story,
dualizing trace gives the pushforward on Jacobians. Together,
pullback and pushforward satisfy the trace identity — the algebraic
backbone of the ``push-then-pull = multiplication-by-degree''
anti-hack theorem.
\end{layman}

\begin{proof}\leanok
Away from the branch locus, define \(\tr_f(\omega)\) by summing the pullbacks
of \(\omega\) along all local inverse branches of \(f\). Symmetric functions
of the inverse branches extend holomorphically across branch points, so the
form extends globally to \(Y\). The constant-map case is treated separately in
the challenge API.
\end{proof}

\begin{proof}[Proof of Theorem~\ref{thm:trace-pullback}]\leanok
Away from the branch locus, choose local inverse branches
\(s_1,\dots,s_d\) of \(f\), where \(d=\deg(f)\). Then
\[
  \tr_f(f^{*}\eta)
  =
  \sum_{i=1}^{d}s_i^{*}f^{*}\eta
  =
  \sum_{i=1}^{d}(f\circ s_i)^{*}\eta
  =
  d\,\eta.
\]
Both sides are holomorphic one-forms on \(Y\), so equality on the complement
of the finite branch locus extends across the branch locus.
\end{proof}

\begin{lemma}[Maps descend to Jacobians]
\label{lem:push-pull-descend}
\uses{def:analytic-jacobian,lem:pullback-forms,lem:trace-forms,lem:pn-r0,lem:pcr-r0,lem:pdp-r0}
Pullback and trace preserve period lattices in the appropriate dual directions,
therefore induce continuous holomorphic group homomorphisms
\[
  f^{*}:\Jac(Y)\to\Jac(X), \qquad f_*:\Jac(X)\to\Jac(Y).
\]
\lean{Jacobian.pullback, Jacobian.pushforward}
\leanok
\end{lemma}

\begin{layman}
Pullback and trace are operations on holomorphic 1-forms; to make them
act on \emph{Jacobians} we have to cross a quotient. The key algebraic
fact is that the dual maps these operations induce on the 1-form duals
send periods to periods — i.e., they preserve the period lattices on
each side. Once that's true, both maps descend to honest homomorphisms
on the lattice quotients, which are the Jacobians. The descended maps
stay holomorphic because they're induced by linear maps between the
covering vector spaces, and a linear map quotients to a smooth one.
Why bother? This is the bridge from ``algebraic operations on 1-forms''
(pullback, trace) to ``geometric operations on Jacobians'' (pullback,
pushforward). Without the descent step, the Jacobian pullback and
pushforward would be undefined; with it, they satisfy the functoriality
and trace identity that the challenge demands.
\end{layman}

\begin{proof}\leanok
The pullback map on forms dualizes to a map \(\HdualX\to\HdualY\), and the
trace map dualizes in the opposite direction. Naturality of integration says
that these dual maps send periods to periods. Therefore they descend to the
quotient tori. Holomorphicity follows because the descended maps are induced
by complex-linear maps between the covering vector spaces.
\end{proof}

\begin{proof}[Proof of Theorem~\ref{thm:push-pull-functoriality}]\leanok
The maps on Jacobians are descended from complex-linear maps between the
universal covering vector spaces, so they are holomorphic group homomorphisms.
The identity laws follow because pullback and trace for the identity map are
the identity on holomorphic one-forms. The composition laws follow from
functoriality of pullback, functoriality of trace for finite holomorphic maps,
and the universal property of quotienting by the period lattices. Equality on
the covering vector spaces descends to equality on the quotient tori.
\end{proof}

\begin{proof}[Proof of Theorem~\ref{thm:pushforward-pullback}]\leanok
Represent \(P\in\Jac(Y)\) by a functional \(\lambda\in\HdualY\). Pullback and
then pushforward acts on the representative by precomposition with
\[
  \HY \xrightarrow{f^{*}} \HX \xrightarrow{\tr_f} \HY .
\]
By Theorem~\ref{thm:trace-pullback}, this composite is multiplication by
\(\deg(f)\). Hence the induced map on the quotient Jacobian is multiplication
by \(\deg(f)\), proving
\[
  f_*(f^{*}P)=\deg(f)\,P.
\]
\end{proof}
